\section{Task 2: Penalty Function $P(\sigma)$}

\subsection{Definition}
For a complete assignment $\sigma$, let
$\ell_{s,k}(\sigma) = \sum_{i:\sigma(t_i)=s} r_k(t_i)$ be the load of
dimension $k$ in slot $s$, and $u_{s,k} = \ell_{s,k}/C_k(s)$ its
utilisation. Let $\mathcal{C}_k = \{s : C_k(s) > 0\}$ be the slots that have
capacity in dimension $k$.
\begin{align*}
P(\sigma) &= P_{\text{base}}(\sigma) + \lambda_R\, R(\sigma) + \lambda_B\, B(\sigma),
  \qquad \lambda_R = \lambda_B = 1\\
P_{\text{base}}(\sigma) &= \sum_{i} w_i\,(\sigma(t_i)+1)\\
R(\sigma) &= \sum_{i:\,u_i > l_i} w_i \left(\frac{\sigma(t_i)-l_i}{u_i-l_i}\right)^{2}
  && \text{(SLA breach risk)}\\
B(\sigma) &= W \cdot \frac{1}{\sum_k |\mathcal{C}_k|}
  \sum_{k=1}^{d} \sum_{s\in\mathcal{C}_k} \left(u_{s,k}-\bar u_k\right)^2,
  \quad \bar u_k = \frac{1}{|\mathcal{C}_k|}\sum_{s\in\mathcal{C}_k} u_{s,k}
  && \text{(load imbalance)}
\end{align*}
Here $[l_i,u_i]$ is the SLA window of $t_i$; the $u_i$ in $R$ is the window
end, not a utilisation.

\subsection{Motivation in the ScoreMe context}
\textbf{SLA breach risk $R$.} A bureau pull placed in the last slot of its
window has no slack left. A single retry against a slow bureau API, or a
delayed upstream OCR job, then breaches the lender SLA. $R$ measures where
each task sits \emph{inside its own window}: 0 at the window start, rising to
$w_i$ at the window end. It is squared because risk is concentrated near the
deadline. Moving a task from the middle of its window to the end costs three
times as much as moving it from the start to the middle.

\textbf{Load imbalance $B$.} If one slot runs at 100\% GPU while the next
runs at 10\%, the OCR cluster has no headroom in the busy slot for retries or
late-arriving jobs, and the idle slot wastes capacity that has already been
paid for. $B$ is the mean squared deviation of utilisation from its
per-dimension mean, scaled by $W$ so that it is in the same units as
$P_{\text{base}}$. A slot with zero capacity in some dimension can hold no
load in it. Its utilisation would be $0/0$, so it is left out of that
dimension. Counting it as 0\% would push work into a slot that cannot take
any.

\subsection{Required properties}
\begin{enumerate}
\item \textbf{Polynomial time.} $P_{\text{base}}$ and $R$ take one pass over
  the tasks, $O(n)$. $B$ builds the load table in $O(nd)$ and then reads it
  in $O(Kd)$. In total, $P$ takes $O(nd + Kd)$ time.
\item \textbf{Monotonically meaningful.} Each term is non-negative and falls
  exactly when the matching operational quantity improves. $P_{\text{base}}$
  falls when weighted delay falls. $R$ falls when tasks move away from their
  deadlines. $B$ is 0 if and only if every dimension is equally utilised in
  every usable slot.
\item \textbf{Non-trivial.} The extension is not a constant, and it is not a
  rescaling of $P_{\text{base}}$. Take two unit tasks $A$ with window $[0,1]$
  and $B$ with window $[0,2]$, and $K=3$. The assignments $(A{\to}1, B{\to}0)$
  and $(A{\to}0, B{\to}1)$ both have $P_{\text{base}} = 3$. But $R = 1$ for
  the first, because $A$ is at its deadline, and $R = 1/4$ for the second. So
  $P$ prefers the second, which $P_{\text{base}}$ cannot tell apart
  (\texttt{test\_extension\_is\_not\_trivial}).
\end{enumerate}

\subsection{Bound used in Task 4}
\begin{lemma}\label{lem:sandwich}
For every feasible $\sigma$:
$P_{\text{base}}(\sigma) \le P(\sigma) \le \beta\, P_{\text{base}}(\sigma)$, where
$\beta = 1 + \lambda_R + \lambda_B/4 = 2.25$.
\end{lemma}
\begin{proof}
$R,B \ge 0$ gives the lower bound. Each ratio in $R$ lies in $[0,1]$, so
$R \le W$. Feasibility (F2) gives $u_{s,k}\in[0,1]$, and the variance of
values in $[0,1]$ is at most $1/4$, so $B \le W/4$. Finally,
$\sigma(t_i)+1\ge 1$ gives $W \le P_{\text{base}}(\sigma)$. Hence
$P \le P_{\text{base}} + \lambda_R W + \lambda_B W/4 \le \beta\,P_{\text{base}}$.
\end{proof}
This lemma lets an approximation ratio for $P_{\text{base}}$ carry over to
$P$ with an extra factor of at most $\beta$ (Task 4). It also makes the
brute-force bound valid for $P$. The bound
$\text{cost} + \sum_{\text{unassigned}} w_i(l_i+1)$ is a lower bound on
$P_{\text{base}}$, and therefore on $P$.

\subsection{Worked example (toy instance)}
For $\sigma = (T1{\to}1, T2{\to}0, T3{\to}0, T4{\to}1, T5{\to}1, T6{\to}2)$,
which is the exact optimum:
$P_{\text{base}} = 33$, and
$R = 5\cdot\tfrac14 + 3\cdot\tfrac19 + 2\cdot\tfrac14 = 2.0833$.
The utilisation vectors over slots $0..3$ are: CPU and RAM
$(0.1875, 1, 0.125, 0)$; GPU $(0,1,0,0)$; Net $(0.833, 0.5, 0.25, 0)$.
They give $B = 19 \cdot 2.3704/16 = 2.8149$, so $P = 37.898$. These values
are asserted in \texttt{tests/test\_penalty.py}.

\subsection{Observation}
On the three small benchmarks (seeds 1, 2, 3), the exact optimum under $P$ is
the \emph{same} assignment as the optimum under $P_{\text{base}}$. The
extension adds 9--15\% to the value but did not change the best schedule.
With $\lambda_R=\lambda_B=1$, delay dominates, and $R$ and $B$ mostly break
ties between near-equal schedules. A platform that treats SLA risk as more
important would raise $\lambda_R$. The bound in Lemma~\ref{lem:sandwich}
then grows to $\beta = 1+\lambda_R+\lambda_B/4$.
